\subsection*{Recognizing and writing two different answers from one passage}
We selected 1,042 pairs of questions about the same traditional Chinese passage for which direct question--answer glyph overlap tied or favoured the other question's answer. Each pair had disjoint acceptable answers. Both writers trained on the same 7,108 separate pairs and replayed the same natural questions. Giving an extra penalty when the sibling answer outranked the requested one increased correct recognition of both answers from 437 to 522 pairs (difference 8.16 percentage points, passage-clustered 95\% interval 4.97--11.39). From a blank start, complete pairs rose from 17 to 21 (difference 0.38 points, interval $-0.59$ to 1.33). The earlier span reader gave 128 complete pairs (Fig.~\ref{fig:same-passage-writer}).

The contrast writer recognized 501 pairs on which it could not independently generate both answers. Among 1,203 incorrect stopped single responses, the model assigned the reference a higher mean log probability than its own response only 30 times. Natural-answer retention also fell. Across all 1,367 human annotations for the 571 DuReader questions, the parent, likelihood and contrast writers matched 35, 23 and 19 complete responses, respectively. None of these exact responses contained a letter, number or Chinese character absent from its input. The original reference alone had credited 32, 22 and 16. Recognizing supplied answers and writing responsive content present different learning problems.
